\bigskip

\begin{enumerate}
	\item 
	\begin{enumerate}
		\item Les vecteurs $\vect{JK}$ et $\vect{JL}$ ont pour coordonnées : \(\overrightarrow{JK}\begin{pmatrix}-1\\2\\0\end{pmatrix}\) \quad \(\overrightarrow{JL}\begin{pmatrix}-4\\-2\\-3\end{pmatrix}\).

Le produit scalaire de  $\vect{JK}$ et $\vect{JL}$ : $\vect{JK}.\vect{JL}=(-1)\times(-4)+2\times (-2)+0\times (-3)=4-4=0$

Les vecteurs $\vect{JK}$ et $\vect{JL}$ sont orthogonaux donc le triangle JKL est rectangle en J.
		\item $\mathcal{A}_{JKL}=\dfrac{JK\times JL}{2}$.

		$JK= \left\|\vect{JK}\right\| =\sqrt{(-1)^2+2^2+0^2}=\sqrt{5}$ et $JL=\left\|\vect{JL}\right\|=\sqrt{(-4)^2+(-2)^2+(-3)^2}=\sqrt{29}$

Donc $\mathcal{A}_{JKL}=\dfrac{\sqrt{5}\times \sqrt{29}}{2}=\dfrac{\sqrt{145}}{2}$ cm\up2.
		\item Le triangle JKL est rectangle en J donc en appliquant la formule de la tangente, on obtient :

		$\tan \left(\widehat{\text{JKL}}\right)=\dfrac{JL}{JK}=\dfrac{\sqrt{29}}{\sqrt{5}}$.

Donc $\widehat{\text{JKL}}=\arctan\left(\dfrac{\sqrt{29}}{\sqrt{5}}\right)\approx 67,5\degres$
	\end{enumerate}

	\item
	\begin{enumerate}
\item Les vecteurs $\vect{JK}$ et $\vect{JL}$ ne sont pas colinéaires : il n'existe pas de réel $k$ tel que $\vect{JK}=k\times \vect{JL}$.

En effet le système $\left\{\begin{array}{l c l} -1&=&k\times (-4)\\2 &=& k\times (-2)\\0&=& k\times (-3) \end{array}\right.$ n'admet pas de solution.

$\left(\vect{JK}~,~\vect{JL}\right)$ est donc une base du plan (JKL).

Calculons :

$\vect{n}.\vect{JK}=6\times(-1)+3\times 2+(-10)\times 0=0$ et $\vect{n}.\vect{JL}=6\times(-4)+3\times (-2)+(-10)\times (-3)=0$.

Le vecteur $\vect{n}$ est donc orthogonal aux $\vect{JK}$ et $\vect{JL}$. C'est donc un vecteur normal au plan (JKL).
		\item Le plan (JKL) a pour équation cartésienne : $ax+by+cz+d=0$, où $(a~;~b~;~c)$ sont les cordonnées d'un vecteur normal au plan. En prenant comme vecteur normal le vecteur $\vect{n}$, on obtient : (JKL) : $6x + 3y - 10z + d = 0$. 

Or $J\in \text{(JKL)}$ donc $12+0-10+d=0$ donc $d=-2$. Donc (JKL) a pour équation 

$6x + 3y - 10z - 2 = 0$.
	\end{enumerate}
	\item
	\begin{enumerate}
		\item La droite $\Delta$ est normale au plan (JKL) donc admet comme vecteur directeur $\vect{n}$, normal à (JKL). elle passe par T$(10~;~9~;~-6)$ donc une représentation paramétrique est :
		
$\left\{\begin{array}{l c l}	x&=&10+6t\\y &=&9+3t\\z&=&-6-10t			\end{array}\right. , t \in \R$
		\item Le projeté orthogonal H de T sur le plan (IJK) est l'unique point d'intersection de $\Delta$ et (IJK). Ses coordonnées sont les solutions du système :

$\left\{\begin{array}{l}	x=10+6t\\y =9+3t\\z=-6-10t\\6x+3y-10z-2=0			\end{array}\right.$

En remplaçant $x$, $y$ et $z$ dans la dernière équation, on obtient :

$6(10+6t)+3(9+3t)-10(-6-10t)-2=0 \iff 145t+145=0 \iff t=-1$

Donc $x=10+6\times(-1)=4$, $y=9+3\times(-1)=6$ et $z=-6-10\times(-1)=4$

H a pour coordonnées $(4~;~6~;~4)$
\item $\mathcal{V}_{JKLT}  = \dfrac13 \mathcal{B} \times h$. La base est le triangle JKL et la hauteur est le segment $[TH]$. D'après les question précédentes :

\(\overrightarrow{TH}\begin{pmatrix}-6\\-3\\10\end{pmatrix}\) donc $TH= \left\|\vect{TH}\right\| = \sqrt{(-6)^2+(-3)^2+10^2}=\sqrt{145}$

Donc $\mathcal{V}_{JKLT}  = \dfrac13 \mathcal{B}_{JKL} \times TH=\dfrac13 \times \dfrac{\sqrt{145}}{2} \times \sqrt{145}= \dfrac{145}{6} $ cm\up3.
	\end{enumerate}

\end{enumerate}
\bigskip
